\documentclass[12pt,a4paper]{report}

\usepackage[utf8]{inputenc}
\usepackage[T1]{fontenc}
\usepackage[french]{babel}
\usepackage{graphicx}
\usepackage{amsmath,amssymb}
\usepackage{listings}
\usepackage{xcolor}
\usepackage{booktabs}
\usepackage{hyperref}
\usepackage{geometry}
\geometry{margin=2.5cm}

\lstset{
  basicstyle=\ttfamily\small,
  breaklines=true,
  frame=single,
  language=C
}

\title{Conception, implémentation et évaluation comparative de structures de données \\
       \large Catalogue d'un commerce en ligne (produits, commandes, clients)}
\author{[Nom 1] \and [Nom 2] \\ Licence 2 Mathématiques-Informatique \\ Université Iba Der Thiam de Thiès}
\date{Année universitaire 2025--2026}

\begin{document}

\maketitle
\tableofcontents

% =====================================================================
\chapter{Introduction}
% Contexte du projet, motivation, domaine applicatif choisi et pourquoi.
% Objectifs du rapport.

% =====================================================================
\chapter{Modélisation des enregistrements}
% Décrire les structs Produit, Client, Commande.
% Justifier chaque champ par rapport aux exigences du sujet (identifiant,
% champs numériques, champs chaînes, champ date, pointeurs inter-types,
% champ tableau/liste interne).
% Inclure un diagramme des relations entre les enregistrements.

% =====================================================================
\chapter{Structures de données implémentées}

\section{Tableau statique de structures}
% Principe, allocation, capacité fixe, accès direct.

\section{Tableau dynamique de pointeurs}
% Principe du redimensionnement géométrique, facteur de croissance documenté.

\section{Liste doublement chaînée}
% Sentinelle / tête-queue, insertion O(1).

\section{[Structure 4 — remplacement de l'ABR]}
% À compléter selon la structure choisie.

\section{[Structure 5 — remplacement du Tas binaire]}
% À compléter selon la structure choisie.

% =====================================================================
\chapter{Étude théorique de complexité}
% Pour chaque couple (structure, opération) : meilleur cas, pire cas, cas
% moyen, en notations O, Omega, Theta, avec justification mathématique
% (récurrences, sommations, théorème maître si pertinent).
% Conclure par un tableau récapitulatif structure x opération.

\begin{table}[h]
\centering
\caption{Synthèse des complexités (à compléter)}
\begin{tabular}{lccccc}
\toprule
Opération & Tab. statique & Tab. dynamique & Liste & [Structure 4] & [Structure 5] \\
\midrule
Insertion       & & & & & \\
Recherche clé    & & & & & \\
Recherche critère secondaire & & & & & \\
Suppression      & & & & & \\
Mise à jour      & & & & & \\
Tri              & & & & & \\
\bottomrule
\end{tabular}
\end{table}

% =====================================================================
\chapter{Protocole expérimental}
% Tailles de jeux de données testées (10^2 à 10^5), distributions utilisées
% (uniforme, gaussienne, triée/inversée), nombre de répétitions, méthode de
% mesure (clock_gettime), conditions matérielles (machine, OS, compilateur).

% =====================================================================
\chapter{Résultats expérimentaux}
% Insérer et commenter les courbes : temps en fonction de n (linéaire et
% log), temps en fonction de la distribution, comparaison théorique vs
% empirique, histogrammes de facteur de ralentissement.

% Exemple d'inclusion de figure :
% \begin{figure}[h]
%   \centering
%   \includegraphics[width=0.8\textwidth]{../../courbes/recherche_temps_n.pdf}
%   \caption{Temps de recherche en fonction de n (échelle log-log)}
% \end{figure}

% =====================================================================
\chapter{Discussion critique}
% Comparer les structures entre elles, expliquer les écarts entre théorie
% et mesures, discuter l'influence de la distribution des données.

% =====================================================================
\chapter{Conclusion}
% Bilan, limites du travail, pistes d'amélioration.

% =====================================================================
\appendix
\chapter{Annexes}
% Extraits de code significatifs, scripts de benchmark, jeux de données.

\end{document}
